
% Question 14
\question
The base radius of the solid right circular cylinder $X$ and the base radius of the solid right circular cone $Y$ are equal. The heights of $X$ and $Y$ are $20\text{ cm}$ and $24\text{ cm}$ respectively. The volume of the solid right circular cone $Z$ is equal to the sum of the volume of $X$ and the volume of $Y$. The base radius of $Z$ is equal to the base diameter of $X$. A craftsman finds that the volume of $Y$ is $800\pi\text{ cm}^3$.
\begin{parts}
\part Find the base radius of $Y$. \rmarks{2}
\part Are $Y$ and $Z$ similar? Explain your answer. \rmarks{3}
\part The craftsman claims that the sum of the curved surface area of $X$ and the curved surface area of $Y$ is greater than the curved surface area of $Z$. Do you agree? Explain your answer. \\*
\null\rmarks{3}
\end{parts}

\begin{solution}
\begin{parts}
\part Let $r\text{ cm}$ be the base radius of $Y$.
\begin{align*}
\frac{24\pi r^2}{3} &= 800\pi \quad &\text{(1M)} \\
r &= 10 \quad &\text{(1A)}
\end{align*}
Thus, the base radius of $Y$ is $10\text{ cm}$.
\part The volume of $Z$ is
\[
\pi(10^2)(20) + 800\pi = 2800\pi\text{ cm}^3 \quad \text{(1M)}
\]
Note that
\[
\left(\frac{\text{The base radius of } Y}{\text{The base radius of } Z}\right)^3 = \left(\frac{1}{2}\right)^3 = \frac{1}{8} \quad \text{and} \quad \frac{\text{The volume of } Y}{\text{The volume of } Z} = \frac{800\pi}{2800\pi} = \frac{2}{7}
\]
Since
\[
\frac{\text{The volume of } Y}{\text{The volume of } Z} \ne \left(\frac{\text{The base radius of } Y}{\text{The base radius of } Z}\right)^3 \quad \text{(1M)}
\]
$Y$ and $Z$ are not similar. \hfill (1A f.t.)
\part The curved surface area of $X$ is
\[
2\pi(10)(20) = 400\pi\text{ cm}^2 \quad \text{(1M)}
\]
The curved surface area of $Y$ is
\[
\pi(10)\left(\sqrt{10^2 + 24^2}\right) = 260\pi\text{ cm}^2 \quad \text{(1M)}
\]
Let $h\text{ cm}$ be the height of $Z$.
\[
\frac{\pi(20^2)(h)}{3} = 2800\pi \implies h = 21
\]
Therefore, the height of $Z$ is $21\text{ cm}$. \\
The curved surface area of $Z$ is
\[
\pi(20)\left(\sqrt{20^2 + 21^2}\right) = 580\pi\text{ cm}^2
\]
The sum of the curved surface area of $X$ and the curved surface area of $Y$ is
\[
400\pi + 260\pi = 660\pi\text{ cm}^2 > 580\pi\text{ cm}^2
\]
Thus, the claim is agreed. \hfill (1A f.t.)
\end{parts}
\end{solution}